% ID: SEL-D0041-OA-E04
% ESTADO: cotejada. Segunda pasada visual y borrador compilado revisados.
% Fuente: Selectividad/matematicas_aplicadas/Ua6MATEMÁTICAS APLICADAS examen.pdf
% SHA-256: b7541cc45d3e96297c29037fe17a562b7daa9cd0ee2a95c2b5e9454ad9b90f14
% Página: 1; opción A; ejercicio original 4.

El capital de las hipotecas constituidas sobre fincas urbanas en Andalucía es una variable aleatoria Normal con desviación típica 10000 \texteuro{}.
\begin{enumerate}[label=\alph*)]
\item \textbf{(2 puntos)} Se toma una muestra aleatoria de 9 hipotecas con los siguientes capitales (en euros):
\[
95000\quad99000\quad105000\quad106000\quad108000\quad111000\quad112000\quad115000\quad120000.
\]
Construya un intervalo de confianza, al 95\%, para el capital medio de dichas hipotecas.
\item \textbf{(0.5 puntos)} \textquestiondown{}Qué número mínimo de hipotecas deberíamos considerar en una muestra para que, con el mismo nivel de confianza, el error máximo en la estimación del capital medio sea de 4000\texteuro{}?
\end{enumerate}
